\subsection{A scalar rate test without computing a slow eigenvector}
\label{sec:rateendpoint}
The endpoint criterion can be evaluated from the generator and one root
of a cubic. This uses the classical polynomial formula for a spectral
projector; the point is an explicit implementation of the preceding
problem-specific criterion, not a new projector construction.

Let $q=\pi(\beta_b)$, $D=\operatorname{diag}(q)$, and put
\[
a=\tfrac14\mathbf1-q,\qquad
d=-\tfrac14(E-\bar E\mathbf1),\qquad
\bar E=\tfrac14\sum_i E_i,\qquad h=D^{-1}a.
\]
Thus $a$ is the infinite-temperature displacement and $d=\pi'(0)$.
Define $\tau=\operatorname{tr}M$,
$\sigma=[(\operatorname{tr}M)^2-\operatorname{tr}(M^2)]/2$ and
$A_j=h^\top M^j d$ for $j=0,1,2$. The scalar polynomial
\[
\mathcal R(x)=A_0x^2+(A_1-\tau A_0)x
                  +A_2-\tau A_1+\sigma A_0
\]
has degree at most two. Let
$p(x)=\det(xI-M)/x$ and let $r$ be its largest root.

\begin{corollary}[Rate evaluation of the triplet endpoint test]
\label{cor:rateendpoint}
For a positive reversible four-state generator satisfying all original
triplet failures, if $r$ is simple, ordinary asymptotic cooling occurs
if and only if $\mathcal R(r)>0$. Equality means absence.
If the largest nonzero root is repeated, cooling is absent by
Theorem~\ref{thm:tripletdegenerate}. Thus the full-premise decision
requires a cubic root and a scalar quadratic evaluation, with no
choice or numerical normalization of a slow eigenvector.
\end{corollary}
\begin{proof}
Write the other two roots as $s,t$, permitting $s=t<r$.
On the invariant zero-sum subspace the spectral projector onto $r$ is
\[
P_r=\frac{(M-sI)(M-tI)}{(r-s)(r-t)}
=\frac{M^2-(\tau-r)M+(\sigma-\tau r+r^2)I}{p'(r)}.
\]
This identity is asserted on that subspace, not on the stationary
direction. Reversibility ensures diagonalizability even if $s=t$.
Since $d$ has zero sum, it follows that
$h^\top P_r d=\mathcal R(r)/p'(r)$.
For any nonzero left slow eigenvector $w$, detailed balance gives
\[
P_r=\frac{Dw w^\top}{w^\top Dw},\qquad
\frac{\mathcal R(r)}{p'(r)}
=\frac{(w^\top a)(w^\top d)}{w^\top Dw}
=\frac{b(0)b'(0)}{w^\top Dw}.
\]
Here $w^\top q=0$, so $b(0)=w^\top a$. Orienting the mode by
the contrast theorem makes $b(0)<0$. Corollary~\ref{cor:tripletendpoint}
therefore makes the product positive exactly when cooling occurs.
Finally $p'(r)=(r-s)(r-t)>0$, proving the sign rule, including equality.
\end{proof}

All quantities except the single cubic root are explicit matrix products
of the original rates, bath weights, and energies. The test still uses
global spectral information through $r$; it is not a new local triplet
test. For rational bath weights and commensurate energies, the common
positive energy unit can be removed and the remaining sign can be
decided by ordered algebraic arithmetic. For arbitrary real energies,
floating evaluation near zero does not settle exact equality. This
corollary concerns existence of an asymptotic effect, not every transient
crossing for a specified distance.
